% !TEX root = userguide.tex

\chapter{How to make a project tree with multiple directories?}
\pgst{tfemchA}
\label{chap:projecttree}

\def\chaptertitle{How to make a project tree with multiple directories?}
\def\headdate{21st November 2023}

The utilities \texttt{getmdefs}, \texttt{fmm} and \texttt{fmme} generate a
macro definition file \texttt{Mdefs.mk} and a \texttt{Makefile} in the current
directory for the files existing in the current directory. Except for small
projects, source files are usually distributed over multiple directories. At
least a separation is usually made between the source files of the modules
and the main program(s). In this section it is shown how to deal with this
situation using \texttt{getmdefs}, \texttt{fmm} (or \texttt{fmme}) and the
additional command \texttt{fmmroot}.

Assume we have the following layout of our project tree (output is from
\texttt{ls -F}):
\begin{quote}
  \texttt{lib/  main/  modfiles/  src/}
\end{quote}
where the directories \texttt{lib} and \texttt{modfiles} are used to store the
project archive and module information files (extension \texttt{.mod}).
The names
of these two directories (\texttt{lib}, \texttt{modfiles}) should be exactly
as specified, whereas other directories can be chosen freely. In this example
the module source files (or other subroutine files) are in \texttt{src} and the
main programs are in \texttt{main}.

The first thing to do after creating the project tree is to generate the
\texttt{Mdefs.mk} in the root project directory by
\begin{quote}
   \texttt{getmdefs -p >Mdefs.mk}
\end{quote}
The option \texttt{-p} will set a macro \texttt{PROJECTDIR} to the project root
directory and a macro \texttt{PROJECTNAME} for the name of the project. The
default name is just \texttt{project}. If you want to use another name, you can
by given the command
\begin{quote}
   \texttt{getmdefs -d >Mdefs.mk}
\end{quote}
instead, which sets \texttt{PROJECTNAME} to the current directory or
\begin{quote}
   \texttt{getmdefs -n \textit{name} >Mdefs.mk}
\end{quote}
which sets \texttt{PROJECTNAME} to the specified name.

Next, create makefiles in all subdirectories containing sources:
\begin{verbatim}
     cd src
     fmm -ls >Makefile
     cd ../main
     fmm -s >Makefile
     cd ..
\end{verbatim}
The option \texttt{-l} will instruct
\texttt{fmm} to create a \texttt{Makefile} that produces an archive library
of the compiled sources and put the \texttt{*.mod} files at a central place.
The option \texttt{-s} will instruct \texttt{fmm} to search for an existing
\texttt{Mdefs.mk} in the tree.

Then, create a `master' \texttt{Makefile} in the project root directory
\begin{quote}
   \texttt{fmmroot >Makefile}
\end{quote}
If there are dependencies between the source directories (there usually are),
you need to edit the master \texttt{Makefile} by hand. In this example the
main programs in \texttt{main} depend on the source files in \texttt{src}.
In order to express that, change the following line 
\begin{quote}
  \texttt{maindir:}
\end{quote}
to
\begin{quote}
   \texttt{maindir: srcdir}
\end{quote}
Now you are all set to build the project by typing
\begin{quote}
   \texttt{make}
\end{quote}
or cleaning the project by
\begin{quote}
   \texttt{make clean}
\end{quote}

Notes:
\begin{itemize}
   \item It is possible to include only a specified set of 
         subdirectories in the master \texttt{Makefile}. See the output of
         \texttt{fmmroot -h}.
   \item In the subdirectories a \texttt{tags} file will be made. By default
only the source files in that directory are included. If you want to include
additional source files, you can add these to the definition of
the \texttt{EXTRATAGSSRC} macro of the \texttt{Mdefs.mk} file in the project 
root directory.
   \item If you type
\begin{quote}
   \texttt{make clean}
\end{quote}
in a subdirectory, all \texttt{.mod} files in \texttt{modfiles} and the
complete library (\texttt{.a}) in \texttt{lib} will be removed. If you don't
like this default, use the \texttt{-n} option:
\begin{verbatim}
     cd src
     fmm -lsn >Makefile
\end{verbatim}
in the source directories. Note, that the \texttt{.mod} files and the library
can now only be cleaned by using the master \texttt{Makefile}.
   \item If the directory containing the main files is somewhere outside of the
         project tree, just copy the \texttt{Mdefs.mk} file to the directory
         (or even better: make a symbolic link) and run \texttt{fmm} in that
         directory.
\end{itemize}


